% Generated draft: requires mathematical and source review.
\documentclass[11pt]{article}
\usepackage[margin=25mm]{geometry}
\usepackage{fontspec}
\setmainfont{Latin Modern Roman}
\usepackage{amsmath,amssymb,mathtools,unicode-math}
\setmathfont{Latin Modern Math}
\usepackage{enumitem,needspace,xurl,hyperref}
\setlength{\parindent}{0pt}
\setlength{\parskip}{5pt}
\setlength{\emergencystretch}{3em}
\newcommand{\N}{\mathbb{N}}
\newcommand{\Z}{\mathbb{Z}}
\newcommand{\Q}{\mathbb{Q}}
\newcommand{\R}{\mathbb{R}}
\newcommand{\C}{\mathbb{C}}
\newcommand{\F}{\mathbb{F}}
\newcommand{\A}{\mathbb{A}}
\newcommand{\PP}{\mathbb{P}}
\newcommand{\E}{\mathbb{E}}
\newcommand{\eps}{\varepsilon}
\newcommand{\dd}{\,\mathrm{d}}
\begin{document}
\section*{UnsolvedMath ID: 30000206}
\textbf{Source problem code:} catalogue entry
\subsection{Definitions and mathematical statement}
The source's Conjecture 1 has an existence-and-normalization part and a compatibility part for the same choice of splitting bundle. Retain its notation \(\operatorname{Hilb}^{(1)}_0\) for the Frobenius-twisted punctual Hilbert scheme and \(V_0\) for a splitting bundle there. In the usual Azumaya-algebra terminology, a splitting bundle for an algebra \(\mathcal A_0\) is a vector bundle equipped with an isomorphism \(\mathcal A_0\simeq\mathcal E nd(V_0)\); the algebra, twisting line bundle and any formal-neighborhood conventions in this problem are not supplied.

Part (a) asks for a choice of \(V_0\) such that, as objects of the intended module category,
\[
 \Gamma\bigl(\operatorname{Hilb}^{(1)}_0,V_0(-1)\bigr)
 \simeq L_{\mathrm{sign}}.
\]
Here \((-1)\) is the source's tensor twist, and \(L_{\mathrm{sign}}\) is its module labelled by the sign representation. This is an ordinary global-sections expression as written, not a replacement by derived global sections or an additional vanishing assertion. The claim includes existence of a normalized splitting bundle, rather than assuming it without comment.

Part (b) asks that, with this choice, the composite derived equivalence referred to as (4) in the source be compatible with the natural \(t\)-structures and the simple-object labels. Write this missing equivalence provisionally as \(\Phi_{V_0}:D^b(\mathcal C)\to D^b(\mathcal D)\). For the intended natural \(t\)-structures, the required compatibility is
\[
 \Phi_{V_0}(D^{\leq0}_{\mathcal C})=D^{\leq0}_{\mathcal D},\qquad
 \Phi_{V_0}(D^{\geq0}_{\mathcal C})=D^{\geq0}_{\mathcal D}.
\]
It should therefore restrict to an equivalence of the hearts, and for every simple \(S_n\)-module \(\tau\) satisfy
\[
 \Phi_{V_0}(L_\tau)\simeq\tau_H.
\]
The original wording indexes \(L_\tau\) by a simple \(S_n\)-module \(\tau\); it does not identify \(L_\tau\) itself with that representation. The labels \(\tau_H\), the abelian categories, the algebra printed imperfectly in the original text, the functors in (4), and the two natural \(t\)-structures are not defined in the supplied record. No replacement equivalence or cohomological shift is guessed.

The Frobenius notation requires a characteristic-\(p\) setup, but the base field, conditions on \(p,n\), parameter \(c\), and any localization hypotheses are missing. The 2026-08-21 literature triage reports the conjecture proved for integral parameters \(c=0,1,2,\ldots\), while not established for all parameters in the cited literature. The admissible all-parameter range and the interpretation of these integral labels in characteristic \(p\) are not specified. The partial-status label and reported direct-source restoration are retained without independent verification.
\subsection{Short English statement}
Choose a splitting bundle with the prescribed sign-labelled global sections, then require its associated derived equivalence to be \(t\)-exact and to match all simple-object labels. Existence, normalization and both compatibility conditions are retained. The omitted algebras, functors, parameter range and characteristic hypotheses prevent a fully self-contained theorem.
\subsection{Sources}
\begin{itemize}
\item Local UnsolvedMath registry, permanent ID 30000206, source code catalogue entry, ``Splitting Bundles for Hilbert-Scheme Equivalences,'' including Conjecture 1(a)--(b), the reference to composite equivalence (4), and the cleaned statement.
\item \emph{Enveloping Algebras and Geometric Representation Theory (2005)}, Oberwolfach Reports, pp.\ 723--777, DOI 10.4171/owr/2005/13: \url{https://doi.org/10.4171/owr/2005/13}.
\item Supplied reference ``Primary characteristic-p paper'': \url{https://arxiv.org/abs/math/0312474}.
\item Literature triage checked 2026-08-21; registry record updated 2026-09-16. The verification field reports restoration of the Frobenius twist and compatibility conditions after checking the report. No external source has been visited or independently checked here.
\end{itemize}
\end{document}
